\documentclass[a4paper,11pt]{article}
\usepackage{plantilla}

\titulo{Apuntes de LaTeX}
\autor{Carlos Sesma}

\begin{document}

\portada
\indice

\section{Estructura de un documento}

Todo documento declara una clase, carga paquetes en el preámbulo y abre el
entorno \texttt{document}. En este repositorio el preámbulo ya está resuelto por
\texttt{plantilla.sty}.

\begin{nota}[Definición]
    El \emph{preámbulo} es todo lo que va entre
    \texttt{\textbackslash documentclass} y
    \texttt{\textbackslash begin\{document\}}. Ahí se carga lo que el documento
    necesita, pero no se escribe contenido.
\end{nota}

\section{Referencias cruzadas}

\texttt{\textbackslash label} marca y \texttt{\textbackslash ref} cita. LaTeX
necesita dos pasadas para resolverlas, por eso se usa latexmk: repite la
compilación hasta que los números dejan de cambiar.

\begin{aviso}[Cuidado]
    Un \texttt{\textbackslash ref} a una etiqueta inexistente no rompe la
    compilación: imprime \texttt{??} y deja solo un aviso en el log. Conviene
    leer los avisos aunque salga el PDF.
\end{aviso}

\section{Código}

El estilo de los listados viene dado por la plantilla.

\begin{lstlisting}[language=Python,caption={Un ejemplo cualquiera}]
def factorial(n):
    return 1 if n <= 1 else n * factorial(n - 1)
\end{lstlisting}

\section{Tablas}

\begin{center}
    \begin{tabular}{lr}
        \textbf{Comando} & \textbf{Pasadas} \\
        \hline
        pdflatex         & 1                \\
        biber            & 1                \\
        pdflatex         & 2                \\
    \end{tabular}
\end{center}

\end{document}
